\documentclass[11pt]{article}

\usepackage[margin=1in]{geometry}
\usepackage{amsmath}
\usepackage{amssymb}
\usepackage{amsthm}
\usepackage{hyperref}
\usepackage{enumitem}

\newtheorem{definition}{Definition}

\title{A Principled Bandwidth for Multitaper Coherence Estimation: \\
Bias from Bessel Zeros, Variance from a Two-Tier Jackknife}
\author{FastMSPEC project design document}
\date{2026-09-03}

\begin{document}
\maketitle

\begin{abstract}
Stage 4's batch run confirmed a deterministic \emph{upper} bound on the multitaper
half-bandwidth $W$ (equivalently the time-bandwidth product $NW$, via $K\approx 2NW$): the
resolution criterion $2W \lesssim \Delta f_{\rm zero} = c_{\min}/(2r)$, where $\Delta f_{\rm
zero}$ is the local spacing between the true coherence's Bessel-function zero-crossings.
What has been missing is a principled \emph{lower} bound --- how small $NW$ can go before the
coherence estimator's own bias/variance behavior, not resolution, degrades picking. This
document works out that lower bound following Haley \& Anitescu's (2017) mean-squared-error
(MSE) framework for multitaper bandwidth selection, adapted in two ways specific to this
project: (i) the bias term is derived in closed form from the known Bessel-zero structure of
the theoretical coherence, rather than estimated by fitting a spline to noisy data, and (ii) the
variance term is decomposed into two tiers --- a taper-level term computable in closed form
(compatible with FastMspec's non-per-taper-materializing architecture) and a record-level term
that separates genuine estimation noise from real day-to-day non-stationarity. Standalone
reference, independent of the conversation history that produced it.
\end{abstract}

\tableofcontents

\section{Motivation}

Ambient-noise interferometry estimates the coherence between two seismic stations, whose real
part follows Aki's (1957) spectral formulation,
\begin{equation}
  \gamma(f) \approx J_0\!\left(\frac{2\pi f r}{c(f)}\right),
  \label{eq:aki}
\end{equation}
matching the convention of \texttt{coherence\_barcode\_design.tex} ($r$: interstation distance,
$c(f)$: phase velocity). The FastMSPEC picker extracts a dispersion curve by tracking this
function's zero-crossings; the multitaper bandwidth $W$ used to estimate $\gamma(f)$ from noisy
data controls, as always in nonparametric spectral estimation, a bias--variance tradeoff.

\textbf{The upper bound is already established.} A bandwidth too wide smears adjacent
zero-crossings together; the resolution criterion
\begin{equation}
  2W \lesssim \Delta f_{\rm zero}(f) = \frac{c(f)}{2r}
  \label{eq:resolution}
\end{equation}
gives $NW_{\rm high}(r) \approx Nc_{\min}/(4r)$ for a window of $N$ samples, where $c_{\min}$ is
the \emph{minimum} phase velocity over the analysis band (not the mean or the range's upper
bound --- the constraint is tightest wherever $c(f)$ is smallest). This was confirmed directly
against Round 1's real data: quartile-4 pairs (780--1030\,km) have $NW_{\rm high}\approx
3.2$--$5.1$, already below the project's fixed operating point of $NW\approx 10.8$, a
quantitative explanation for their observed near-zero convergence rate.

\textbf{The lower bound is not.} Today's $NW$ values (13--15 for MspecBestK, 80 for Mspec,
FastMspec's much smaller effective taper count) are hand-picked, not derived. This document
addresses that gap.

\section{The MSE Framework (Haley \& Anitescu, 2017)}

Haley \& Anitescu \cite{haley2017} give a systematic method for multitaper bandwidth selection:
minimize the estimated MSE of the log spectrum,
\begin{equation}
  \hat\Lambda(f; N, W) = \frac{W^4}{36}\left[\frac{S''(f)}{\bar S(f)}\right]^2 +
  \hat\sigma_J^2(f),
  \label{eq:mse}
\end{equation}
summed over a frequency grid, where the first term is squared bias (from a second-order Taylor
expansion of the true spectrum, Section~\ref{sec:bias}) and $\hat\sigma_J^2(f)$ is a jackknife
variance estimate (Section~\ref{sec:variance}). The optimal $NW$ minimizes the sum. Two features
of their method do not transfer directly to this project and are addressed below: (i) their bias
term is estimated from a spline fit to the \emph{observed}, noisy multitaper spectrum itself
--- appropriate when nothing more is known about the true spectrum, but this project has a much
stronger prior, the known Bessel-zero structure of Equation~\ref{eq:aki}; and (ii) their variance
theory is for a single-record auto-spectrum, whereas this project's coherence estimate is a
cross-spectral quantity additionally stacked over many independent day/window records
(\texttt{coh\_num}, often hundreds to $\sim$1600) --- a second averaging layer their paper does
not address.

\section{Bias: A Closed Form from the Bessel-Zero Structure}
\label{sec:bias}

Haley \& Anitescu's local bias term (their Eq.~10, here Eq.~\ref{eq:mse}'s bias piece) is
\begin{equation}
  B_W(f) \approx \frac{W^2}{6}S''(f),
  \label{eq:bias-generic}
\end{equation}
the residual left by a symmetric smoothing kernel acting on the true spectrum's local curvature.
We show this applies \emph{exactly} at a true zero-crossing $f_n$ of Equation~\ref{eq:aki}, not
merely by analogy, and that $\gamma''(f_n)$ is computable in closed form.

\subsection{The bias formula holds exactly at $f_n$}

Write the true coherence near its zero at $f_n$ as $\gamma(f) \approx \gamma'(f_n)(f-f_n) +
\tfrac12\gamma''(f_n)(f-f_n)^2$. Smoothing with a symmetric kernel $K_W$ (the multitaper spectral
window, Eq.~7 of the Haley--Anitescu supplement):
\begin{equation}
  \bar\gamma(f_n) = \int K_W(f_n-\nu)\gamma(\nu)\,d\nu \approx
  \gamma'(f_n)\underbrace{\int K_W(u)u\,du}_{=0} +
  \tfrac12\gamma''(f_n)\int K_W(u)u^2\,du.
\end{equation}
The first-derivative (slope) term vanishes identically for any symmetric kernel --- an odd
integrand over a symmetric window. What remains is exactly Eq.~\ref{eq:bias-generic}, evaluated
at $f_n$: the smoothed estimate does not sit at zero, it sits at a spurious offset set purely by
local curvature. This spurious residual is what a picker can mistake for a shifted or missing
crossing.

\subsection{Closed-form $\gamma''(f_n)$}

With $\gamma(f) = J_0(x(f))$, $x(f) = 2\pi f r/c(f)$, and the Bessel recurrence $J_1'(x) = J_0(x)
- J_1(x)/x$: at a zero $x_n$ of $J_0$ (where $J_0(x_n)=0$),
\begin{equation}
  \gamma''(f_n) = J_1(x_n)\left[\frac{(x'(f_n))^2}{x_n} - x''(f_n)\right].
  \label{eq:gamma-pp}
\end{equation}
Every term is computable without new data: $x_n$ is the $n$-th zero of $J_0$ (exact,
\texttt{scipy.special.jn\_zeros}); $J_1(x_n)$ is a standard special-function evaluation, never
zero at $x_n$ by Bessel zero-interlacing; $x'(f_n)$, $x''(f_n)$ are derivatives of $2\pi f
r/c(f)$, using the exact station distance $r$ and $c(f)$'s local derivatives from the reference
or picked curve already used for $c_{\min}$ in the Round 2 subset-selection script
(\texttt{c\_ref(f) + best\_delta\_km\_s}).

\textbf{Assumption, stated explicitly}: this requires $c(f)$ to be locally well-behaved enough
near each $f_n$ for its derivatives to be meaningfully estimated from the picked curve --- weaker
than a globally constant $c$ (already rejected for the $c_{\min}$ correction), but not
guaranteed, especially in a rapidly-dispersive band. Worth checking empirically, not assumed.

\subsection{Payoff}

This replaces Haley \& Anitescu's self-referential spline-on-noisy-data bias estimate with a
closed-form bias term derived from known physics, evaluated exactly where it matters (the
zero-crossings themselves). No fitting step, no dependence on the very estimate whose quality
it's meant to characterize.

\section{Variance: A Two-Tier Jackknife}
\label{sec:variance}

\subsection{Taper level (compatible with FastMspec's architecture)}

The delete-one jackknife variance of the log spectrum over $K$ tapers (Haley--Anitescu
supplement Eq.~23) needs each individual per-taper spectrum --- straightforward for classical
Mspec (which already materializes per-taper spectra, chunked over $K$), but in tension with
FastMspec's deliberately non-per-taper-materializing architecture (Section~\ref{sec:fastmspec}).
The supplement's Eq.~24 resolves this: the \emph{expected value} of the jackknife variance is a
closed form in $K$ alone (via the trigamma function $\psi'$),
\begin{equation}
  E\{\hat\sigma_J^2(f)\} = \frac{(K-1)^3}{K\binom{K-1}{2}}\left[\frac{2}{(K-2)^2} +
  \frac12\left(\psi'\!\left(\frac{K-1}{2}\right) - \psi'\!\left(\frac{K-2}{2}\right)\right)\right],
  \label{eq:jk-expected}
\end{equation}
requiring only $K$ (or FastMspec's effective transition-region taper count, denoted $r_{\rm
trans}$ to avoid collision with the distance symbol $r$ above), not per-taper recomputation.

\textbf{Update (2026-09-03), a real correction found by testing Eq.~\ref{eq:jk-expected} against
a literal Eq.~23 jackknife computed from real per-taper cross-spectra (SKRH--BAND pair, Mspec's
classical materialized-per-taper approach): Eq.~\ref{eq:jk-expected} is not distribution-free,
and it does not transfer cleanly to cross-spectra.} Its derivation assumes each per-taper
quantity is $\chi^2(2)$-distributed --- true for an auto-spectrum (the squared magnitude of one
complex Gaussian) but not for a cross-spectrum per-taper term (the product of \emph{two}
correlated complex Gaussians), whose distribution near a coherence null is qualitatively
different --- the same $\sim 1/K$ near-null instability documented in Walden (2000). Binning
the literal jackknife variance by per-taper magnitude (a proxy for proximity to a null) on real
data shows the lowest-magnitude quartile has jackknife variance roughly $20$--$30\times$ the
middle quartiles, and the literal-vs-closed-form ratio grows with $K$ (1.9$\times$ at $K=5$,
8.3$\times$ at $K=33$) rather than converging --- Eq.~\ref{eq:jk-expected} systematically
underestimates variance exactly where the picker's crossing decisions are made.

\textbf{Revised recommendation}: use Eq.~\ref{eq:jk-expected} away from candidate zero-crossings,
where its assumptions are closer to valid and FastMspec's memory-decoupling argument stays
intact, but fall back to a literal per-crossing jackknife (or a genuinely cross-spectral closed
form --- Walden (2000)'s Wishart-based multivariate treatment \cite{walden2000} is the right
place to check, not Haley \& Anitescu's univariate theory) in the neighborhood of each candidate
crossing specifically. This is a bounded cost (a handful of crossings per pair, not the whole
band), not a full retreat from Section~\ref{sec:fastmspec}'s argument -- but it does mean
Eq.~\ref{eq:jk-expected} alone is not sufficient at the frequencies that matter most, and should
not be presented as though it were.

\subsection{Record level (a real complication, not glossed over)}

Haley \& Anitescu's theory is single-record. This project stacks \texttt{coh\_num} independent
day/window traces on top of the $K$-taper average --- a second layer their paper never addresses.
Following the precedent in Park \& Levin (2016) \cite{parklevin2016}, who jackknife over
\emph{events} in a bin-average (not tapers) for exactly this multi-record situation (``a
bin-average with 10 events leads to 10 different RF estimates,'' via the delete-one method of
Efron 1982), decompose the per-trace estimates $X_i(f)$ (log- or Fisher-transformed) as
\begin{equation}
  X_i(f) = \mu(f) + \delta_i(f) + \varepsilon_i(f),
\end{equation}
where $\delta_i$ is genuine day-to-day physical variation (non-stationarity: storms, source
shifts) and $\varepsilon_i$ is that day's own $K$-taper estimation noise. Then
\begin{equation}
  \widehat{\mathrm{Var}}(\delta; f) \approx \mathrm{SampleVar}_{i}(X_i(f)) -
  E\{\hat\sigma_J^2(f)\},
  \label{eq:excess-variance}
\end{equation}
using Eq.~\ref{eq:jk-expected} as the estimation-noise baseline (cheap, FastMspec-compatible) and
the plain sample variance across traces as the empirical total (Haley \& Anitescu's jackknife
motivation, statistical efficiency for a \emph{small} number of independent samples, applies most
naturally at the $K$-taper tier; at \texttt{coh\_num} scale, already in the hundreds, plain sample
variance is simpler and no less appropriate).

\textbf{This is a diagnostic, not just a variance estimate.} Computing
Equation~\ref{eq:excess-variance} as a function of frequency identifies where the coherence
estimate is unstable for a \emph{physical} reason (real non-stationarity) rather than an
estimator-tunable one (insufficient tapering). At frequencies where $\widehat{\mathrm{Var}}
(\delta; f)$ dominates, no choice of $NW$ will stabilize picking --- the MSE-minimization
machinery of Section~\ref{sec:mse-recipe} is the wrong tool there, and that should be stated
plainly in any picking-reliability discussion rather than attributed to a suboptimal bandwidth.

\section{Why FastMspec's Architecture Survives This Analysis}
\label{sec:fastmspec}

FastMspec (Karnik, Romberg \& Davenport's algorithm \cite{karnik2022}, translated in
\texttt{fast\_multitaper.py}) decomposes the multitaper average into an $O(N\log N)$ bulk
sinc-convolution term plus a small correction over only the DPSS eigenvalue transition-region
tapers, whose count $r_{\rm trans}$ stays bounded ($11$--$24$) across $NW=5$--$1600$ while $K$
grows linearly the whole way. This means (i) the memory argument for FastMspec over classical
Mspec/MspecBestK holds regardless of where the bandwidth search in this document lands, and (ii)
the taper-level variance term above (Eq.~\ref{eq:jk-expected}) never needs FastMspec's fused
computation broken open --- $K$ alone (or $r_{\rm trans}$, if the transition-region-only variance
is preferred) is sufficient.

\section{Putting It Together}
\label{sec:mse-recipe}

The full recipe, combining Sections~\ref{sec:bias}--\ref{sec:variance} with the already-derived
upper bound (Equation~\ref{eq:resolution}):
\begin{enumerate}[nosep]
  \item Compute $NW_{\rm high}(r)$ per pair from its own $c_{\min}$ (Round 2 subset-selection
    script).
  \item At each candidate $NW$ in the sweep, compute the closed-form bias
    (Equation~\ref{eq:gamma-pp}) at the theoretical zero-crossing frequencies, and the closed-form
    taper-level variance (Equation~\ref{eq:jk-expected}).
  \item Minimize $\hat\Lambda(NW) = \mathrm{bias}^2 + \mathrm{variance}$, restricted to $NW \le
    NW_{\rm high}(r)$ (the resolution bound is a hard ceiling on catastrophic failure, not part
    of the smooth tradeoff; the true optimum can sit well inside it, as in Haley \&
    Anitescu's own Lorentzian example, $NW=14.5$ against a much larger ceiling).
  \item Separately compute the excess-variance diagnostic (Equation~\ref{eq:excess-variance}) to
    flag frequencies where bandwidth tuning cannot help at all.
\end{enumerate}
Round 2's NW sweep (60-pair stratified subset, 5 fractions of $NW_{\rm high}(r)$ each) is the
mechanism for computing this empirically once real per-fraction convergence/quality data exists
--- this document gives the theory the sweep is meant to test, not a substitute for running it.

\section{Open Questions}

\begin{itemize}[nosep]
  \item Whether $c(f)$'s local derivatives, needed for Equation~\ref{eq:gamma-pp}, are reliably
    estimable from the picked/reference curve in rapidly-dispersive bands, or need their own
    smoothing/regularization.
  \item Whether $\delta_i$ and $\varepsilon_i$ (Section~\ref{sec:variance}) are actually
    independent, as assumed in Equation~\ref{eq:excess-variance} --- untested against real data.
  \item Whether the closed-form bias (Section~\ref{sec:bias}) should use the fundamental-mode
    Bessel structure alone, or needs correction for multi-mode contamination, matching the same
    caveat already carried by the Love-wave Bessel-model simplification in
    \texttt{coherence\_barcode\_design.tex}.
  \item \emph{Tested, not fully resolved}: the extent of the near-null breakdown documented in
    Section~\ref{sec:variance}'s update -- confirmed present and large on one real pair (SKRH--
    BAND), not yet checked across the dataset's real distance/coherence range, nor against a
    genuinely cross-spectral closed-form alternative.
\end{itemize}

\begin{thebibliography}{9}

\bibitem{aki1957}
Aki, K. (1957). Space and time spectra of stationary stochastic waves, with special reference to
microtremors. \textit{Bulletin of the Earthquake Research Institute}, 35, 415--456.

\bibitem{haley2017}
Haley, C.\ L., \& Anitescu, M.\ (2017). Optimal bandwidth for multitaper spectrum estimation.
\textit{IEEE Signal Processing Letters}, 24(11), 1696--1700.
\url{https://doi.org/10.1109/LSP.2017.2719943}

\bibitem{karnik2022}
Karnik, S., Romberg, J., \& Davenport, M.\ A.\ (2022). Multitaper spectral estimation revisited:
fast algorithms and generalized approaches. Preprint. (FastMSPEC project reference:
\texttt{Multitaper\_Revisited\_Karnik.pdf}.)

\bibitem{parklevin2016}
Park, J., \& Levin, V.\ (2016). Statistics of multiple-taper receiver functions. \textit{Geophysical
Journal International}, 207(1), 512--527. \url{https://doi.org/10.1093/gji/ggw291}

\bibitem{walden2000}
Walden, A.\ T.\ (2000). A unified view of multitaper multivariate spectral estimation.
\textit{Biometrika}, 87(4), 767--788. \url{https://doi.org/10.1093/biomet/87.4.767}

\bibitem{ekstrom2009}
Ekstr\"om, G., Abers, G.\ A., \& Webb, S.\ C.\ (2009). Determination of surface-wave phase
velocities across USArray from noise and Aki's spectral formulation. \textit{Geophysical Research
Letters}, 36(18), L18301. \url{https://doi.org/10.1029/2009GL039131}

\end{thebibliography}

\end{document}
